% GENERATED FILE. Edit canonical lessons, then run npm run generate:books.
% canonical-source-hash: fnv1a64:b7b8d65304d06c4b
% canonical-lessons: KA-C154-yaru, KA-C154-endigu, KA-C154-elliyu, KA-C154-yaradaru, KA-C154-enadaru

\chapter{Nobody, Never, Nowhere, Anyone, Anything}
\label{ch:ka-nobody-never-nowhere-anyone-anything}
\hlchaptermodality{154}

\begin{chapteropening}
\textbf{By the end of this chapter:} \emph{I can say what is not so, and ask yes-or-no questions.}

Complete the last lesson of chapter 154: I can say what is not so, and ask yes-or-no questions.
\end{chapteropening}

\section[\texorpdfstring{\kn{ಯಾರೂ} (yārū)}{ಯಾರೂ (yārū)}]{\kn{ಯಾರೂ} (yārū) — nobody (with a negative verb)}

\label{lesson:KA-C154-yaru}



\begin{quote}
Before the new one: say the Kannada for to pull, then the Kannada for to cry.
\end{quote}

\subsection*{You'll want to know: \kn{ಯಾರೂ}}
\textbf{\kn{ಯಾರೂ}} — \emph{yārū} — \textquotedblleft{}nobody (with a negative verb)\textquotedblright{}. Kannada negates a verb by joining \textbf{\kn{ಇಲ್ಲ}} (\emph{illa}) onto it: \textbf{\kn{ಯಾರೂ} \kn{ಬರಲಿಲ್ಲ}} — \emph{yārū baralilla} — \textquotedblleft{}Nobody came.\textquotedblright{} \textbf{\kn{ಇಲ್ಲ}} also means \textquotedblleft{}is not there\textquotedblright{}; to say something is not a thing, it uses \textbf{\kn{ಅಲ್ಲ}} (\emph{alla}): \textbf{\kn{ಇದು} \kn{ನನ್ನ} \kn{ಮನೆ} \kn{ಅಲ್ಲ}} — \emph{idu nanna mane alla} — \textquotedblleft{}This is not my house.\textquotedblright{}

\begin{grammarlens}[title={What you've built}]
One more word for this chapter.
\end{grammarlens}

\subsection*{Your turn}
\begin{itemize}
\raggedright
  \item \emph{Say it:} \emph{yārū}
  \item \emph{Say it:} \emph{yārū}, once more
  \item \emph{Say it:} say \emph{yārū}
  \item \emph{Recall:} say the Kannada for to pull, then the Kannada for to cry, then say \emph{yārū} again
\end{itemize}

\begin{tcolorbox}[breakable,colback=teal!4,colframe=teal!35!black,title={Before you move on}]
What does \kn{ಯಾರೂ} mean? (\textquotedblleft{}Nobody (with a negative verb)\textquotedblright{}.) Say it once more.
\end{tcolorbox}
\section[\texorpdfstring{\kn{ಎಂದಿಗೂ} (endigū)}{ಎಂದಿಗೂ (endigū)}]{\kn{ಎಂದಿಗೂ} (endigū) — never (with a negative verb)}

\label{lesson:KA-C154-endigu}



\begin{quote}
Before the new one: say the Kannada for to cry, then the Kannada for nobody.
\end{quote}

\subsection*{You'll want to know: \kn{ಎಂದಿಗೂ}}
\textbf{\kn{ಎಂದಿಗೂ}} — \emph{endigū} — \textquotedblleft{}never (with a negative verb)\textquotedblright{}. \textbf{\kn{ನಾನು} \kn{ಎಂದಿಗೂ} \kn{ಮಾಂಸ} \kn{ತಿನ್ನುವುದಿಲ್ಲ}} — \emph{nānu endigū māṃsa tinnuvudilla} — \textquotedblleft{}I never eat meat.\textquotedblright{}

\begin{grammarlens}[title={What you've built}]
One more word for this chapter.
\end{grammarlens}

\subsection*{Your turn}
\begin{itemize}
\raggedright
  \item \emph{Say it:} \emph{endigū}
  \item \emph{Say it:} \emph{endigū}, once more
  \item \emph{Say it:} say \emph{endigū}
  \item \emph{Recall:} say the Kannada for to cry, then the Kannada for nobody, then say \emph{endigū} again
\end{itemize}

\begin{tcolorbox}[breakable,colback=teal!4,colframe=teal!35!black,title={Before you move on}]
What does \kn{ಎಂದಿಗೂ} mean? (\textquotedblleft{}Never (with a negative verb)\textquotedblright{}.) Say it once more.
\end{tcolorbox}
\section[\texorpdfstring{\kn{ಎಲ್ಲಿಯೂ} (elliyū)}{ಎಲ್ಲಿಯೂ (elliyū)}]{\kn{ಎಲ್ಲಿಯೂ} (elliyū) — nowhere, not anywhere (with a negative verb)}

\label{lesson:KA-C154-elliyu}



\begin{quote}
Before the new one: say the Kannada for nobody, then the Kannada for never.
\end{quote}

\subsection*{You'll want to know: \kn{ಎಲ್ಲಿಯೂ}}
\textbf{\kn{ಎಲ್ಲಿಯೂ}} — \emph{elliyū} — \textquotedblleft{}nowhere, not anywhere (with a negative verb)\textquotedblright{}. \textbf{\kn{ಅವನು} \kn{ಎಲ್ಲಿಯೂ} \kn{ಹೋಗಲಿಲ್ಲ}} — \emph{avanu elliyū hōgalilla} — \textquotedblleft{}He didn't go anywhere.\textquotedblright{}

\begin{grammarlens}[title={What you've built}]
One more word for this chapter.
\end{grammarlens}

\subsection*{Your turn}
\begin{itemize}
\raggedright
  \item \emph{Say it:} \emph{elliyū}
  \item \emph{Say it:} \emph{elliyū}, once more
  \item \emph{Say it:} say \emph{elliyū}
  \item \emph{Recall:} say the Kannada for nobody, then the Kannada for never, then say \emph{elliyū} again
\end{itemize}

\begin{tcolorbox}[breakable,colback=teal!4,colframe=teal!35!black,title={Before you move on}]
What does \kn{ಎಲ್ಲಿಯೂ} mean? (\textquotedblleft{}Nowhere, not anywhere (with a negative verb)\textquotedblright{}.) Say it once more.
\end{tcolorbox}
\section[\texorpdfstring{\kn{ಯಾರಾದರೂ} (yārādarū)}{ಯಾರಾದರೂ (yārādarū)}]{\kn{ಯಾರಾದರೂ} (yārādarū) — anyone, someone}

\label{lesson:KA-C154-yaradaru}



\begin{quote}
Before the new one: say the Kannada for never, then the Kannada for nowhere.
\end{quote}

\subsection*{You'll want to know: \kn{ಯಾರಾದರೂ}}
\textbf{\kn{ಯಾರಾದರೂ}} — \emph{yārādarū} — \textquotedblleft{}anyone, someone\textquotedblright{}. A yes-or-no question changes the last vowel of the sentence's final word to \textbf{-\kn{ಆ}} (\emph{-ā}): \textbf{\kn{ಇದ್ದಾರೆ}} becomes \textbf{\kn{ಯಾರಾದರೂ} \kn{ಇದ್ದಾರಾ}?} — \emph{yārādarū iddārā?} — \textquotedblleft{}Is anyone there?\textquotedblright{}

\begin{grammarlens}[title={What you've built}]
One more word for this chapter.
\end{grammarlens}

\subsection*{Your turn}
\begin{itemize}
\raggedright
  \item \emph{Say it:} \emph{yārādarū}
  \item \emph{Say it:} \emph{yārādarū}, once more
  \item \emph{Say it:} say \emph{yārādarū}
  \item \emph{Recall:} say the Kannada for never, then the Kannada for nowhere, then say \emph{yārādarū} again
\end{itemize}

\begin{tcolorbox}[breakable,colback=teal!4,colframe=teal!35!black,title={Before you move on}]
What does \kn{ಯಾರಾದರೂ} mean? (\textquotedblleft{}Anyone, someone\textquotedblright{}.) Say it once more.
\end{tcolorbox}
\section[\texorpdfstring{\kn{ಏನಾದರೂ} (ēnādarū)}{ಏನಾದರೂ (ēnādarū)}]{\kn{ಏನಾದರೂ} (ēnādarū) — anything, something}

\label{lesson:KA-C154-enadaru}



\begin{quote}
Before the new one: say the Kannada for nowhere, then the Kannada for anyone.
\end{quote}

\subsection*{You'll want to know: \kn{ಏನಾದರೂ}}
\textbf{\kn{ಏನಾದರೂ}} — \emph{ēnādarū} — \textquotedblleft{}anything, something\textquotedblright{}. \textbf{\kn{ನಿಮಗೆ} \kn{ಏನಾದರೂ} \kn{ಬೇಕಾ}?} — \emph{nimage ēnādarū bēkā?} — \textquotedblleft{}Do you need anything?\textquotedblright{}

\begin{grammarlens}[title={What you've built}]
One more word for this chapter.
\end{grammarlens}

\subsection*{Your turn}
\begin{itemize}
\raggedright
  \item \emph{Say it:} \emph{ēnādarū}
  \item \emph{Say it:} \emph{ēnādarū}, once more
  \item \emph{Say it:} say \emph{ēnādarū}
  \item \emph{Recall:} say the Kannada for nowhere, then the Kannada for anyone, then say \emph{ēnādarū} again
\end{itemize}

\begin{tcolorbox}[breakable,colback=teal!4,colframe=teal!35!black,title={Before you move on}]
What does \kn{ಏನಾದರೂ} mean? (\textquotedblleft{}Anything, something\textquotedblright{}.) Say it once more.
\end{tcolorbox}
